\subsection{What the compiler outputs for the Standard Model and our reality}
\label{subsec:smmap}
Every entry below is an \emph{output} of $(\cthree,\gcar)$ through the $D_5\times A_3$ compiler; none is
an independent input. ``Reads'' names the lattice object the observable comes from.

\begin{center}\small
\begin{tabularx}{\linewidth}{@{}p{0.30\linewidth}p{0.40\linewidth}p{0.13\linewidth}c@{}}
\toprule
\textbf{SM / cosmology observable} & \textbf{Compiler output} & \textbf{Reads} & \textbf{St.} \\
\midrule
Number of generations & $\Nfam=3=\operatorname{rank}A_3=\dim H^1(\mathbb P^1\setminus\mu_4)$ & $A_3$ & \mE \\
Gauge algebra dimension & $\dimg=12=|R(A_3)|=\gcd(6r,2h)$ & $A_3$ & \mE \\
One generation ($16$, incl.\ $\nu^c$) & $S^+=D_5$ half-spinor $=\Lambda^{0,2,4}E$ & $D_5$ & \mE \\
Hypercharge spectrum & all $Y$ from $y_\pm=\{\tfrac12,-\tfrac13\}$; $\operatorname{Tr}X^2=120$ & $D_5$ root data & \mE \\
Abelian index $b_1$ & $10b_1=1+|\mu_4||E(K_5)|=1+4\cdot10=41$ & $\mu_4\times K_5$ & \mE \\
Fermion family occupancy & $\Oadm=\Nfam\dim S^+=(\Nfam{+}N_\Phi)|R(A_3)|=48$ & $A_3$ & \mE \\
Fine structure $\ainv=137.0359992168$ & root of $F_{U(1)}(\alpha)=0$ (explicit, unique) & EM closure & \mE \\
Cabibbo angle & $\lC=\sqrt{\useed(1-\useed)}$, $\useed=\tfrac43\cthree+48\cthree^4$ & seed & \mE \\
CKM CP phase & $\delta_{\rm CKM}=\tfrac{2\pi}{|R^+(A_3)|}+\Nfam\useed(1-\useed)=\tfrac\pi3+3\lC^2$ & $A_3^+$, seed & \mE/\mC \\
PMNS angles/phase & $\sin^2\theta_{12}=\tfrac13-\tfrac\useed2$, $\sin^2\theta_{13}=e^{-5/6}\useed$, $\delta^\nu=\tfrac{4\pi}{3}$ & $A_3^+$, seed & \mE/\mC \\
Charged-lepton hierarchy & $\det M_\ell\propto\useed^{\,h/\Nfam}=\useed^{10}$; $(5,3,2)$ & decuple & \mE/\mC \\
Mass-hierarchy transport & $\sum L=40=|R(D_5)|$; entries on $C_6=|R^+(A_3)|$ & $D_5\times A_3$ & \mE \\
Strong CP & $\theta_{\rm eff}=0$ (determinant erasure, given RP; document~3, Part~III) & --- & \mE/\mC \\
Weak CP & holonomy: $\delta_{\rm CKM}-\tfrac\pi3=\Nfam\,\useed(1-\useed)$ & $A_3^+$ & \mC \\
Gravity tree normaliser & $\xi_{\rm tree}=\tfrac{\dimg}{\dim S^+}=\tfrac34=q(A_3)$ & $A_3$ glue & \mE \\
Cosmological constant & $\rho_\Lambda/\Mbar^4=\tfrac{3}{4\pi^2}e^{2\dph}[v_{\rm geo}/(5\betarad^2\Mbar)]^{10}$ & decuple & \mE \\
Hubble anchor & $H_0\sqrt{\Omega_\Lambda}=\Mbar e^{-\ainv}/(2\pi)$ & action & \mE \\
Baryon density & $\omega_b=\lC/\AL=\lC/10$ & decuple & \mE/\mC \\
Inflation (scalaron) & $M_{\rm scal}^2/\Mbar^2=\cthree^{\,\Oadm-10b_1}=\cthree^{7}$ & deficit $r-1$ & \mE \\
Inflation amplitude & $A_s=N_\star^2\cthree^7/24\pi^2\approx2.2\times10^{-9}$ & deficit & \mC/\mO \\
\bottomrule
\end{tabularx}
\end{center}
\markerkey

\noindent\emph{Reading.} The discrete sector --- generations, gauge dimension, family content,
hypercharge, $b_1$, occupancy --- is now \emph{group-theoretic output} of $D_5\times A_3$. The
continuous sector --- $\alpha^{-1}$, $\lC$, the mixings, $\rho_\Lambda$, $H_0$, inflation --- is the
seam seed $\cthree$ and the action exponentials read on the \emph{same} carrier graphs ($K_5$ pairs,
the $C_6$ hexagon $=|R^+(A_3)|$, the decuple $\AL=\binom{\gcar}{2}=10$). What used to be a list of
separate hits is, structurally, one machine: $D_5$ carries the family and the flavor budget, $A_3$
carries the generations and the gauge dimension, $\mu_4$ glues them (and reappears as the $41$ chain,
the discriminant norms, and the CP phase step $2\pi/6$), and $\cthree$ dresses everything.

% =====================================================================
\section{Grammar I --- the carrier grammar ($3+2$)}
% =====================================================================
\noindent\emph{Machine checks for this section:} \veri{v2\_\allowbreak carrier\_\allowbreak pascal.py}


From $\gcar=5$ and the split $(b,s)=(3,2)$ the carrier data follow:
\begin{equation}
\begin{aligned}
  &\gamma_{\mathrm{car}}=\frac{\gcar}{\gcar+1}=\frac56,\quad
  \Nfam=\frac{\gcar+1}{2}=3,\quad
  \dim S^+=2^{\gcar-1}=16,\\[2pt]
  &\Oadm=\Nfam\,2^{\gcar-1}=48,\qquad
  b_1=\frac{\gcar\,2^{\gcar-2}+1}{10}=\frac{41}{10}.
\end{aligned}
\end{equation}

\subsection{Dual-root dictionary (everything from two charges)}

\par\smallskip\noindent\textbf{Two carrier roots generate the family and the transport cusps \mE.}
With $y_+=\tfrac12$ and $y_-=-\tfrac13$ (the determinant-normalised carrier eigenvalues), the full
one-family hypercharge spectrum and the flavor-transport cusps are sums of the two roots:
\[
  \begin{aligned}
  Y(Q_L)&=y_++y_-=\tfrac16, & Y(u^c)&=2y_-=-\tfrac23, & Y(d^c)&=-y_-=\tfrac13,\\
  Y(L_L)&=-y_+=-\tfrac12, & Y(e^c)&=2y_+=1, & Y(\nu^c)&=0,
  \end{aligned}
\]
\[
  \{1,\tfrac23,\tfrac13\}=\{\,2y_+,\ -2y_-,\ 2(y_++y_-)\,\},\qquad
  \gamma_{\mathrm{car}}=\operatorname{Tr}_E Y^2=3y_-^2+2y_+^2=\tfrac56,
\]
\[
  \phibase=\frac{1-\gamma_{\mathrm{car}}}{\pi}=\frac{1}{6\pi}.
\]
All residuals are exactly zero. The Standard-Model charge table is not assumed; it is built by
addition from $\{1/2,-1/3\}$.
\par\smallskip

\begin{keybox}[The $X=6Y$ quadratic and the hypercharge Pascal-sign duality \mE]
Clearing denominators with $X=6Y$ turns the two base charges into the integer roots $x_+=3$, $x_-=-2$ of
\[
  \boxed{\ x^2-x-6=0\ }:\qquad
  x_++x_-=1=N_\Phi,\quad x_+-x_-=5=\gcar,
\]
\[
  -x_+x_-=6=|R^+(A_3)|,\qquad \sqrt{\Delta}=5,
\]
so the Higgs singleton, the carrier rank and the $A_3$ hexagon all fall out of one quadratic. Counting
the signs of $X$ over a generation ($X=1^{\times6},2^{\times3},6^{\times1},(-4)^{\times3},(-3)^{\times2},0^{\times1}$)
gives the \emph{Pascal triple} again:
\[
  \boxed{\ (\#X{=}0,\,\#X{<}0,\,\#X{>}0)=(1,5,10)=\tbinom50,\tbinom51,\tbinom52\ },
\]
\[
  10{-}5{=}\gcar,\quad 5{-}1{=}|\mu_4|,\quad 10{-}1{=}\operatorname{tr}R,
\]
with $\nu^c\leftrightarrow\binom50$ the neutral singleton --- the right-handed neutrino is the Pascal zero
with a job. The moment ladder adds $\operatorname{Tr}X^2=120=5!$ and
$\operatorname{Tr}X^4/\operatorname{Tr}X^2=2280/120=19$ (an $E_8$ exponent; audit). \mE
\end{keybox}

\subsection{Carrier arithmetic (proved identities)}

The discrete coefficients are not isolated numbers; they are exact carrier traces:
\begin{align}
  10\,b_1 &= 41 = \gcar^2+(\gcar-1)^2 = \Oadm\,\gamma_{\mathrm{car}}+1
      = 1+h_{E_8}\tfrac{\dim S^+}{\dimg} = 1+30\cdot\tfrac{16}{12} = 40+1, &&\mE\\
  I_1^{\mathrm{EW}} &= \frac{41}{48}=\operatorname{Var}_{S^+}(Y)\cdot b_1=\frac{5}{24}\cdot\frac{41}{10}
      = \gamma_{\mathrm{car}}+\frac{1}{\Oadm}, &&\mE\\
  \operatorname{Tr}_{S^+}(X^2) &= 120 = 5!\quad(X=6Y), \qquad
  \dtwo = \tfrac54\dtop^2 = 2880\,\cthree^8 = 4!\cdot 5!\,\cthree^8 = 4!\,\operatorname{Tr}_{S^+}(X^2)\,\cthree^8, &&\mE\\
  \frac54 &= \underbrace{\tfrac32\cdot\tfrac56}_{B\gamma}
      = \underbrace{\tfrac{\Dstart}{\Oadm}}_{60/48}
      = \underbrace{\tfrac{\gcar}{\gcar-1}}_{5/4}
      = \frac{\dtwo}{\dtop^2}. &&\mE
\end{align}
So: $b_1$ is a Pythagorean carrier identity ($5^2+4^2$); equivalently the anchor-product form
$10b_1=p_1p_3+(\gcar-p_1)=4\cdot10+(5-4)=41$ reads ``all four glue corners times all ten carrier pairs,
plus the one leftover Higgs slot''; the electroweak exponent $41/48$ is the
hypercharge variance times the $U(1)$ budget; the second seam defect $2880\,\cthree^8$ is a product
of factorial traces $4!\cdot5!$; and the universal compression $5/4=\gcar/(\gcar-1)$ recurs whenever
the five-slot carrier projects onto four effective directions. As a minor curiosity, the SU(5)-tree
weak mixing angle rewrites as $\sin^2\theta_W^{\mathrm{tree}}=\tfrac38=\Nfam/(\gcar+\Nfam)=\Nfam/\operatorname{rank}(E_8)$
(the value is the standard one; the rewrite is a coincidence to note, not a claim). \mC

% =====================================================================
\section{Grammar II --- the seed grammar (one decoder $\useed$)}
% =====================================================================
\noindent\emph{Machine checks for this section:} \veri{v106\_\allowbreak review\_\allowbreak validation.py}


The seed is, exactly, a pure function of $\cthree$:
\begin{equation}
  \boxed{\ \useed=\varphi^{\mathrm{ret}}_0=\frac43\,\cthree+\Oadm\,\cthree^4
  =\underbrace{\frac{1}{6\pi}}_{\phibase}+\underbrace{\frac{3}{256\pi^4}}_{\dtop}=0.0531719521768\ }\qquad\mE.
\end{equation}
Its two coefficients are themselves carrier data: $\tfrac43=\dim S^+/\dimg=16/12=8/(\gcar+1)$ (so
$\phibase=1/((\gcar+1)\pi)$) and $48=\Oadm$. The seed is thus a pure function of $\cthree$ and the
carrier counts.

\par\smallskip\noindent\textbf{One seed, four readouts (and their inverses) \mE.}
\[
  \lC^2=\useed(1-\useed),\quad
  \betarad=\frac{\useed}{4\pi},\quad
  \Omega_b=\Bigl(1-\frac{1}{4\pi}\Bigr)\useed,\quad
  \sin^2\theta_{13}=e^{-5/6}\,\useed,
\]
\[
\begin{gathered}
  \useed=\frac{1-\sqrt{1-4\lC^2}}{2}=4\pi\,\betarad=\frac{\Omega_b}{1-1/(4\pi)}=e^{5/6}\sin^2\theta_{13}\\[2pt]
  =\sum_{n\ge0}C_n\,\lC^{2n+2}=\lC^2+\lC^4+2\lC^6+5\lC^8+\cdots
\end{gathered}
\]
\par\smallskip

\noindent One knob $\useed$ turns four dials simultaneously: Cabibbo mixing (the Bernoulli width
$\sqrt{\useed(1-\useed)}$), cosmic birefringence, the baryon fraction, and the reactor angle. Turning
one forces the others; there is no free per-observable tuning. The exponent $5/6$ is the carrier
trace $\gamma_{\mathrm{car}}=\operatorname{Tr}_E Y^2$, \emph{not} the Euler--Mascheroni constant
$0.5772\ldots$ (audit \mO).

% =====================================================================
\section{Grammar III --- the action grammar (exponentials of $\ainv$)}
% =====================================================================

\noindent\textbf{The relation in plain words.} The fine-structure constant is not just ``a number near
$137$'': it is the \emph{one exponential engine} behind three wildly separated scales. Take the same
$\ainv=137.036$ and divide the carrier into it: $\tfrac15$ of it sets the \emph{electroweak} scale,
$1\times$ it sets the \emph{Hubble} scale, $2\times$ it sets the \emph{cosmological constant}. The three
action charges are the Pascal row $1:5:10$ of the five-slot carrier, and the Hubble rung is just the
square root of the vacuum energy. So one constant fixes the electroweak hierarchy, dark energy and the
expansion rate at once.

\begin{center}
\begin{tikzpicture}[font=\small,
  eng/.style={rectangle,rounded corners=3pt,draw=tfgold!80!black,fill=tfgold!16,line width=1pt,align=center,inner sep=6pt},
  sc/.style={rectangle,rounded corners=2pt,draw=tfgreen!70!black,fill=tfgreen!8,align=center,inner sep=4pt,text width=46mm},
  ar/.style={-{Stealth[length=2.6mm]},thick,tfgray}]
  \node[eng] (e) at (0,0) {\textbf{one engine}\\$\ainv=137.0359992$};
  \node[sc] (ew) at (6.2,2.0) {\textbf{electroweak} \;($q{=}\tfrac1{\gcar}$)\\$v_{\rm EW}/\Mbar\sim e^{-\ainv/5}$};
  \node[sc] (h)  at (6.2,0.0) {\textbf{Hubble} \;($q{=}1$)\\$H_0/\Mbar\sim e^{-\ainv}=\sqrt{\rho_\Lambda}$};
  \node[sc] (l)  at (6.2,-2.0){\textbf{cosm.\ constant} \;($q{=}2$)\\$\rho_\Lambda/\Mbar^4\sim e^{-2\ainv}$};
  \draw[ar] (e) to[bend left=12] node[above,font=\scriptsize]{$\div\,5$} (ew);
  \draw[ar] (e) -- node[above,font=\scriptsize]{$\times\,1$} (h);
  \draw[ar] (e) to[bend right=12] node[below,font=\scriptsize]{$\times\,2$} (l);
  \node[font=\scriptsize,tfgray,align=center] at (3.0,-3.25)
    {action charges $1:\gcar:2\gcar=1:5:10=\binom50:\binom51:\binom52$ (Pascal row of $K_5$)};
\end{tikzpicture}
\end{center}

Read every dimensionless ratio through its \emph{action} and ask for the action charge:
\begin{equation}
  \boxed{\ x=R_x\,e^{-A_x},\qquad A_x=q_x\,\ainv+\Delta_x,\qquad q_x\in\Bigl\{\tfrac{1}{\gcar},\,1,\,2\Bigr\}.\ }
\end{equation}
The robust ladder and its residues are
\begin{center}\small
\begin{tabular}{@{}lllll@{}}
\toprule
sector & action $A_x$ & charge $q_x$ & residue $R_x$ & value of $-\ln x$ \\
\midrule
electroweak $v_{\mathrm{geo}}/\Mbar$ & $(\ainv+\dph)/\gcar$ & $1/\gcar$ & $\gcar\betarad^2$ & $36.85$ \\
Hubble $H_0/\Mbar$ & $\ainv$ & $1$ & $1/(2\pi\sqrt{\Omega_\Lambda})$ & $138.68$ \\
cosm.\ const.\ $\rho_\Lambda/\Mbar^4$ & $2\ainv$ & $2$ & $3/(4\pi^2)$ & $276.65$ \\
\bottomrule
\end{tabular}
\end{center}
\begin{equation}
  \AEW=\frac{\ainv+\dph}{\gcar}=27.5259,\qquad \AL=2\ainv=274.072,\qquad
  \boxed{\ \AEW{:}\AH{:}\AL=1{:}5{:}10\ }.
\end{equation}
The leading ratio is exact: $\AL/(\ainv/\gcar)=2\gcar=10$ (with the transport correction
$\AL/\AEW=9.957$). The claim is that the \emph{action charges}, not the raw logarithms, form the
ladder.

\begin{figure}[ht]
\centering
\includegraphics[width=0.82\linewidth]{figures/action_tower.pdf}
\caption{Action tower (data plot, \texttt{verification/make\_figures.py}). The three action charges
$A_x=-\ln x$ for the electroweak, Hubble and cosmological-constant ratios fall on the line
$A=(\ainv/5)\binom5k$: one engine $\ainv$ generates all three scales with the Pascal exponents $1:5:10$.}
\label{fig:actiontower}
\end{figure}

